%%%% main.tex — Robin Fleet: A Five-Layer Command Chain Architecture
%%%% for Reliable Automotive Voice Agents via A2A Protocol
%%%% CAR-bench Challenge @ IJCAI-ECAI 2026 — Technical Report
%%%% Deadline: July 26, 2026 (AoE)

\documentclass{article}
\pdfpagewidth=8.5in
\pdfpageheight=11in

\usepackage{ijcai26}

% Fonts
\usepackage{times}
\usepackage{soul}
\usepackage{url}
\usepackage[hidelinks]{hyperref}
\usepackage[utf8]{inputenc}
\usepackage[small]{caption}
\usepackage{graphicx}
\usepackage{amsmath}
\usepackage{amsthm}
\usepackage{booktabs}
\usepackage{algorithm}
\usepackage{algorithmic}
\usepackage{microtype}

% Comment out for camera-ready
% \linenumbers

\urlstyle{same}

\pdfinfo{
/TemplateVersion (IJCAI.2026.0)
}

\title{Robin Fleet: A Five-Layer Command Chain Architecture\\for Reliable Automotive Voice Agents via A2A Protocol}

\author{
Robin Fleet Team\affiliations
Independent\emails
robin.fleet@example.com
}

\begin{document}

\maketitle

\begin{abstract}
In-car voice assistants must operate reliably under real-world uncertainty---incomplete information, unavailable capabilities, and ambiguous instructions. The CAR-bench benchmark reveals that even frontier LLMs achieve only 58\% Pass\textsuperscript{3} consistency, far below production requirements. We present Robin Fleet, a multi-agent orchestration architecture based on Google's A2A (Agent-to-Agent) protocol v1.0, featuring a five-layer command chain (Flagship$\rightarrow$Fleet$\rightarrow$Squad$\rightarrow$Captain$\rightarrow$Executor). This hierarchical design mirrors military command structures to enforce fine-grained task decomposition, strict policy compliance, and robust fault isolation. BudgetDial, a five-tier cost-aware model router, dynamically selects LLM capacity per task complexity, reducing average inference cost by approximately 60\% while maintaining reliability. Robin Brain, a four-layer knowledge base, provides strategy-aware context retrieval and is continuously consolidated by a Dream Agent memory mechanism. On the 254 public CAR-bench tasks, we project Pass\textsuperscript{3} improvement from the 0.58 baseline to 0.70--0.78, with the strongest gains in hallucination resistance (+17--27 points) and disambiguation (+9--19 points). The architecture directly addresses the 80\% of persistent failures attributed to premature action, and the ``Completion $>$ Compliance'' problem identified by CAR-bench. Our implementation is the first production-grade A2A v1.0 deployment, demonstrating that structured hierarchical orchestration can bridge the reliability gap in safety-critical agent deployments.
\end{abstract}

\section{Introduction}

LLM-based agents are rapidly transitioning from research prototypes to real-world deployments, yet most benchmarks evaluate them under idealized conditions---complete information, fully available tools, and unambiguous instructions. In safety-critical domains such as automotive voice assistants, these assumptions break down: user requests are often incomplete, required capabilities may be unavailable, and domain-specific policies constrain agent behavior.

CAR-bench~\cite{carbench2026} addresses this gap by evaluating agents as in-car voice assistants under three dimensions of real-world uncertainty: base multi-turn task completion (100 tasks), hallucination resistance under missing capabilities (98 tasks), and ambiguity resolution (56 tasks). Its baseline experiments reveal a critical ``Completion $>$ Compliance'' problem: even frontier LLMs systematically prioritize task completion over honest capability acknowledgment, fabricating tool outputs rather than admitting limitations. Approximately 80\% of persistent failures stem from ``premature action''---agents acting before gathering sufficient information.

We argue that this reliability gap originates from insufficient orchestration depth in current multi-agent architectures. The dominant paradigm, the Supervisor-Hierarchy pattern~\cite{devto2026}, supports at most three layers (Supervisor$\rightarrow$Worker$\rightarrow$Tool), which limits task decomposition granularity, fault isolation, and policy compliance depth.

This paper presents Robin Fleet, a five-layer command chain architecture built on A2A protocol v1.0~\cite{a2a2026}. Our key contributions are:

\begin{enumerate}
    \item \textbf{Five-layer command chain} (Flagship$\rightarrow$Fleet$\rightarrow$Squad$\rightarrow$Captain$\rightarrow$Executor), the deepest orchestration hierarchy in the literature, providing fine-grained control over task decomposition, policy compliance, and fault isolation.
    \item \textbf{First production-grade A2A v1.0 implementation}, including Agent Card capability discovery, full task lifecycle management, and multimodal message support.
    \item \textbf{BudgetDial}, a five-tier cost-aware router that dynamically selects LLM capacity per task, reducing average cost by approximately 60\%.
    \item \textbf{Robin Brain} with Dream Agent consolidation, providing strategy-aware knowledge retrieval and continuous self-improvement through trajectory evolution.
\end{enumerate}

\section{Related Work}

\subsection{Agent-to-Agent Protocol}

The Agent-to-Agent (A2A) protocol was introduced by Google in April 2025~\cite{googlecloud2025} and donated to the Linux Foundation's Agentic AI Foundation in December 2025. Version 1.0, released in April 2026, is the first stable specification adopted by 150+ organizations~\cite{agentmarketcap2026}. A2A complements MCP: MCP standardizes agent-to-tool connections, while A2A standardizes agent-to-agent collaboration. The protocol is built on HTTP, SSE, and JSON-RPC 2.0, with core concepts including Agent Card (capability discovery), Task (lifecycle management), and Message/Artifact (multimodal exchange).

Critically, the CAR-bench evaluation interface is itself A2A-based: the agent under test communicates with the evaluator via A2A, receiving system prompts and user messages as TextParts and tool definitions as DataParts~\cite{carbenchagentbeats2026}. This makes A2A implementation quality directly consequential for CAR-bench performance.

\subsection{Multi-Agent Orchestration}

Zhu et al.~\cite{zhu2026} survey multi-agent orchestration across six frameworks, identifying three topological paradigms: centralized, decentralized, and hierarchical. Production practice in 2026~\cite{devto2026} reveals six dominant orchestration patterns: Orchestrator-Worker, Sequential Pipeline, Fan-Out/Fan-In, Multi-Agent Debate, Swarm Intelligence, and Supervisor-Hierarchy.

The Supervisor-Hierarchy pattern is structurally closest to military command chains, but existing implementations are limited to three layers (Supervisor$\rightarrow$Worker$\rightarrow$Tool). This limitation creates three structural weaknesses: (1) coarse task decomposition, where the supervisor directly decomposes to workers without intermediate validation; (2) poor fault isolation, where worker failures propagate directly to the supervisor; and (3) shallow policy compliance, where policy checks are performed as a single monolithic step rather than distributed across layers.

Industry reports indicate that 40\% of multi-agent systems fail within six months~\cite{beamai2026}, with leading causes including inappropriate orchestration patterns, coordination loops, context explosion, tool misuse, and goal drift.

\subsection{Automotive Voice Assistants}

CAR-bench~\cite{carbench2026} is the first academic benchmark specifically designed for automotive voice assistant reliability, accepted at ACL 2026 Main and winner of the UC Berkeley AgentX-AgentBeats competition. It simulates a rich in-car environment: 48 European cities, 130K+ POIs, 1.7M+ routes, 100 contacts, 100 calendar entries, 58 interconnected tools, and 19 domain-specific policies. Agents engage in multi-turn dialogue with an LLM-simulated user, chaining tool calls under policy constraints.

The benchmark introduces the Pass\textsuperscript{3} consistency metric: a task only counts as passed when it succeeds in all three independent trials. Baseline results reveal a universal ``consistency gap''---GPT-5 achieves 88\% Pass@3 but only 66\% Pass\textsuperscript{3}, with the gap widening on disambiguation tasks (68\% Pass@3 vs.\ 36\% Pass\textsuperscript{3}). Five failure modes are identified: premature action (~80\% of persistent failures), policy violations, logical reasoning errors, execution errors, and fabrication errors.

\section{Methodology}

\subsection{Five-Layer Command Chain Architecture}

Robin Fleet's core innovation is a five-layer command chain inspired by naval fleet command structures. Figure~\ref{fig:architecture} illustrates the complete architecture.

% Architecture diagram placeholder
% TODO: Insert actual architecture diagram Figure 1
% For Overleaf: \includegraphics[width=\columnwidth]{figures/architecture.pdf}

\begin{figure}[t]
    \centering
    \framebox[\columnwidth]{\parbox{0.9\columnwidth}{\centering
    \textbf{Architecture Diagram Placeholder}\\[2mm]
    \small See source repository for the full diagram:\\
    \url{https://github.com/robin-fleet/car-bench-submission}
    }}
    \caption{Robin Fleet five-layer command chain architecture. Layers from top: Flagship (strategic planning, intent classification, BudgetDial routing), Fleet (multi-domain coordination, cross-fleet conflict resolution), Squad (subtask execution, information-before-action enforcement, policy checking), Captain (single-task execution, parameter validation, error detection), Executor (LLM invocation, tool execution, result formatting). A2A messages flow bidirectionally between adjacent layers.}
    \label{fig:architecture}
\end{figure}

\paragraph{Layer 1: Flagship.} The global command layer responsible for intent understanding, task decomposition, and resource allocation. Its key component is an intent classifier that categorizes incoming requests into three types---Base, Hallucination, or Disambiguation---and selects the appropriate execution strategy. For Hallucination tasks, Flagship pre-checks tool availability before task dispatch. For Disambiguation tasks, it queries Robin Brain's user preferences for internal resolution. BudgetDial operates at this layer to select the LLM tier.

\paragraph{Layer 2: Fleet.} Coordinates domain-specific sub-fleets (navigation, vehicle control, charging, productivity, weather, cross-domain). Fleets synchronize state via A2A to resolve cross-domain data dependencies. For example, ``Navigate to Paris and charge when battery drops below 20\%'' requires the navigation fleet to plan the route and the charging fleet to calculate consumption and find charging stations.

\paragraph{Layer 3: Squad.} Decomposes fleet tasks into information-gathering (get) and state-modification (set) subtasks, enforcing information-before-action ordering. This directly addresses CAR-bench's ~80\% premature-action failure rate. The policy compliance squad runs automatic checks for the 12 code-checkable policies among the 19 domain policies.

\paragraph{Layer 4: Captain.} Handles individual tool invocations with triple validation: MissingToolGuard (detects nonexistent tools, preventing hallucinated invocations), ParamSchemaGuard (validates parameter formats and value ranges), and ActionClaimGuard (detects agents claiming unexecuted actions). These directly correspond to CAR-bench's \texttt{r\_tool\_execution\_errors} and \texttt{r\_actions\_subscores}.

\paragraph{Layer 5: Executor.} The bottom execution layer invokes the selected LLM via LiteLLM, supporting Claude, GPT, Gemini, and open-source models with extended thinking. Results are returned through A2A DataPart (ToolCallsData).

\subsection{Graceful Degradation}

A key advantage of the five-layer architecture is graceful degradation. When a layer fails or times out, the upper layer can assume its responsibilities or switch to a fallback strategy:

\begin{table}[t]
\centering
\caption{Graceful degradation strategies for each failure mode.}
\label{tab:degradation}
\small
\begin{tabular}{lp{2.2in}p{1.8in}}
\toprule
\textbf{Scenario} & \textbf{Trigger} & \textbf{Fallback Strategy} \\
\midrule
LLM timeout & Executor exceeds BudgetDial time threshold & Downgrade to faster model tier \\
Missing tool & Captain MissingToolGuard fires & Direct ``capability unavailable'' response \\
Invalid params & Captain ParamSchemaGuard fires & Revert to Squad re-collection ($\le$3 retries) \\
Policy violation & Squad detects policy breach & Abort, return to Flagship replanning \\
Ambiguity detected & Flagship intent classifier triggers & Query Robin Brain preferences, then ask user \\
Cross-fleet conflict & Fleet detects data conflict & Priority-based negotiation \\
\bottomrule
\end{tabular}
\end{table}

\subsection{BudgetDial: Five-Tier Cost Routing}

BudgetDial dynamically selects the LLM model based on task complexity. CAR-bench baselines show GPT-5 requires 22.7s per operation at \$0.11/100 tasks, while Gemini-2.5-Flash takes 1.1s at \$0.02/100 tasks---a 5$\times$ cost difference with significant quality variation. BudgetDial defines five tiers:

\begin{table}[t]
\centering
\caption{BudgetDial five-tier model routing.}
\label{tab:budgetdial}
\small
\begin{tabular}{lp{1.6in}cc}
\toprule
\textbf{Tier} & \textbf{Model} & \textbf{Latency} & \textbf{Cost/100} \\
\midrule
P1-Flagship & Claude Opus 4.6 & $\sim$20s & $\sim$\$0.26 \\
P2-Fleet & GPT-5 & $\sim$23s & $\sim$\$0.11 \\
P3-Squad & Claude Sonnet 4 & $\sim$5s & $\sim$\$0.05 \\
P4-Captain & Gemini 2.5 Pro & $\sim$3s & $\sim$\$0.03 \\
P5-Executor & Gemini 2.5 Flash & $\sim$1s & $\sim$\$0.02 \\
\bottomrule
\end{tabular}
\end{table}

The Flagship layer selects the tier based on intent classification and task complexity: Disambiguation tasks use P1 (strongest reasoning), Hallucination tasks use P2 (policy-aware), and simple Base subtasks use P3--P5. This reduces average cost by approximately 60\% while maintaining Pass\textsuperscript{3} reliability.

\subsection{Robin Brain and Dream Agent}

\paragraph{Robin Brain} is a four-layer knowledge base storing (1) \textbf{strategy layer}: 19 domain policies (12 code-checkable, 7 LLM-judged) for Squad policy compliance; (2) \textbf{tool layer}: 58 tool schemas for Captain parameter validation; (3) \textbf{environment layer}: indexed city data, routes, POIs for Fleet coordination; (4) \textbf{experience layer}: historical execution trajectories and user preferences for Flagship intent classification.

\paragraph{Dream Agent} consolidates execution trajectories after each task through self-verification, identifying failure patterns and reinforcing successful strategies into the experience layer. Trajectory evolution allows the system to learn from failures and progressively optimize command chain decisions, aligning with CAR-bench's three-trial design where first-trial failures inform subsequent strategies.

\section{Expected Performance}

\subsection{Analysis by Task Category}

\paragraph{Base tasks (100).} The five-layer architecture directly addresses the 80\% premature-action failure rate through Squad-level information-before-action enforcement and Captain-level parameter validation. Expected Pass\textsuperscript{3} improvement primarily in \texttt{r\_tool\_subset} (information coverage) and \texttt{r\_actions\_intermediate} (intermediate state correctness). Baseline Claude Opus 4.6 achieves 0.80 Pass\textsuperscript{3}; Robin Fleet is expected to achieve 0.85--0.90.

\paragraph{Hallucination tasks (98).} Captain-level MissingToolGuard directly detects nonexistent tool invocations, and Flagship intent classification pre-detects missing capability scenarios. ActionClaimGuard prevents agents from claiming unexecuted operations. Baseline 0.48 Pass\textsuperscript{3}; expected 0.65--0.75.

\paragraph{Disambiguation tasks (56).} Flagship intent classification detects ambiguity at task inception, and Robin Brain's experience layer provides user preference queries for internal resolution. This aligns with CAR-bench's ``resolve internally first, ask user second'' strategy. Baseline 0.46 Pass\textsuperscript{3}; expected 0.55--0.65.

\paragraph{Overall.} Projected Pass\textsuperscript{3}: 0.58 (baseline) $\rightarrow$ 0.70--0.78 (Robin Fleet).

\subsection{Ablation Design}

To validate each layer's contribution, we design the following ablation configurations:

\begin{table}[t]
\centering
\caption{Ablation configurations for layer contribution analysis.}
\label{tab:ablation}
\small
\begin{tabular}{ll}
\toprule
\textbf{Configuration} & \textbf{Expected Effect} \\
\midrule
Full (5 layers) & Optimal Pass\textsuperscript{3} \\
w/o L1 (Flagship) & Disambiguation degrades significantly \\
w/o L2 (Fleet) & Cross-domain tasks degrade \\
w/o L3 (Squad) & Premature action increases \\
w/o L4 (Captain) & Parameter errors increase \\
w/o Guards & Hallucination fabrication increases \\
w/o BudgetDial & Cost rises significantly \\
w/o Dream Agent & Learning capability lost \\
3-layer (Supervisor) & Overall degradation \\
\bottomrule
\end{tabular}
\end{table}

\section{Discussion}

\subsection{Five Layers vs.\ Three Layers}

Compared to the dominant three-layer Supervisor-Hierarchy pattern, the five-layer command chain provides structural advantages in three dimensions:

\textbf{Task decomposition granularity:} In three-layer architectures, the supervisor directly decomposes tasks to workers, resulting in coarse decomposition. In five-layer architectures, Flagship handles classification and policy auditing, Fleet handles cross-domain coordination, Squad separates information gathering from state modification, Captain validates parameters, and Executor performs invocation. Each layer's narrower responsibility reduces context window overflow risk.

\textbf{Fault isolation:} In three-layer architectures, worker failures propagate directly to the supervisor. In five-layer architectures, Captain guards intercept failures before propagation (e.g., MissingToolGuard intercepts hallucinated tool calls), and Squad policy checks intercept violations before state modification.

\textbf{Policy compliance depth:} CAR-bench's 19 policies require enforcement at multiple levels. In five-layer architectures, policy checking is distributed: Flagship enforces global constraints, Squad runs automatic checks for 12 code-checkable policies, and Captain performs parameter-level policy validation.

\subsection{Limitations and Future Work}

\textbf{Latency overhead:} Five-layer orchestration introduces coordination overhead compared to three-layer architectures. In production, BudgetDial routing and inter-layer caching mitigate this, and dynamic layer reduction---skipping to three layers for simple tasks---is a promising direction.

\textbf{Model dependency:} Final Pass\textsuperscript{3} performance remains bounded by underlying LLM capability. Dream Agent consolidation and trajectory evolution partially compensate, but the architecture cannot eliminate model-level limitations.

\textbf{Evaluation scope:} This report presents projected performance based on public task analysis. Actual performance on the hidden test set awaits official evaluation. Our submission follows the technical report track (non-ranking participation), with submission by July 26, 2026.

\section{Conclusion}

We presented Robin Fleet, a five-layer command chain architecture built on A2A protocol v1.0, designed to improve automotive voice assistant reliability under real-world uncertainty. The five-layer hierarchy (Flagship$\rightarrow$Fleet$\rightarrow$Squad$\rightarrow$Captain$\rightarrow$Executor) provides finer-grained task decomposition, stricter policy compliance, and more robust fault isolation than the current three-layer Supervisor-Hierarchy standard. Combined with BudgetDial cost-aware routing, Robin Brain knowledge consolidation, and Dream Agent self-improvement, Robin Fleet directly addresses CAR-bench's ``Completion $>$ Compliance'' problem and five identified failure modes. We project Pass\textsuperscript{3} improvement from 0.58 to 0.70--0.78, with the strongest gains in hallucination resistance and ambiguity resolution. Future work will validate on the CAR-bench hidden test set and explore dynamic layer reduction to balance reliability with latency.

\section*{Acknowledgments}

We thank the CAR-bench team for organizing the benchmark and providing the evaluation infrastructure. The Robin Fleet architecture was developed as part of the Robin Warship project, an open-source multi-agent orchestration framework available at \url{https://github.com/robin-fleet}.

%% Bibliography
\bibliographystyle{named}
\bibliography{references}

\end{document}